% !TeX program = lualatex
%
% Exemplo da classe crpsp-timbrado: papel timbrado do CRP SP, com os cinco
% níveis de título. A arte (crpsp-timbrado-logo.png) não vem no repositório;
% sem ela o cabeçalho cai para o nome do Conselho em texto.
%
\DocumentMetadata{
  lang        = pt-BR,
  pdfversion  = 2.0,
  pdfstandard = ua-2,
  tagging     = on,
}
\documentclass{crpsp-timbrado}

\hypersetup{
  pdftitle  = {Exemplo de documento em papel timbrado},
  pdfauthor = {CRP SP},
}

\begin{document}

\noindent São Paulo, 15 de outubro de 2026.

\bigskip
\noindent Ofício n.\,º 000/2026

\section{Objeto}

Este exemplo mostra a classe \texttt{crpsp-timbrado} com os cinco níveis de
título, todos em linha própria, e o texto corrido em Atkinson Hyperlegible Next.
Títulos e rodapé são compostos em NewJune.

\subsection{Cabeçalho e rodapé}

O logo fica no cabeçalho e o endereço no rodapé. O endereço é texto, e não
imagem como no modelo do designer: dá para selecionar e pesquisar. Cabeçalho e
rodapé são artefatos de página e ficam fora da ordem de leitura; por isso o
remetente deve estar identificado no corpo do texto, na assinatura.

\subsubsection{Páginas seguintes}

A partir da segunda página aparece a indicação “Página X de Y” à direita do
logo, o que ajuda quem navega sem ver a folha inteira.

\paragraph{Quarto nível}

Texto sob o título de nível 4. Na classe \texttt{article} padrão este título
ficaria embutido na linha; aqui ele é um título como os outros.

\subparagraph{Quinto nível}

Texto sob o título de nível 5.

\section{Encaminhamentos}

São os seguintes os encaminhamentos:

\begin{enumerate}
  \item avaliar a minuta do documento;
  \item solicitar parecer da Gerência Jurídica;
  \item publicar a versão final em formato acessível.
\end{enumerate}

\subsection{Prazos}

\begin{tabular}{@{}ll@{}}
  \textbf{Etapa}  & \textbf{Prazo}    \\
  Avaliação       & 30 de outubro     \\
  Parecer         & 13 de novembro    \\
  Publicação      & 27 de novembro    \\
\end{tabular}

\section{Texto longo}

Parágrafo de enchimento para forçar uma segunda página e exercitar a indicação
de número de página no cabeçalho. A psicóloga inscrita no Conselho Regional
exerce função pública e está sujeita à fiscalização ética; o órgão zela pela
qualidade técnica, pela justiça e pela inclusão das pessoas.

\count255=0
\loop\ifnum\count255<12
  \advance\count255 by 1
  Três princípios sustentam a atuação: o compromisso com os direitos humanos, a
  defesa da ciência e a construção coletiva. Não há exercício profissional sem
  responsabilidade, nem avaliação séria sem informação clara.\par
\repeat

\bigskip
\noindent Atenciosamente,

\bigskip\bigskip
\noindent Nome da pessoa signatária\\
Cargo\\
Conselho Regional de Psicologia de São Paulo (CRP SP)

\end{document}
